\section*{Capítulo \ref{ch:b2:muscle-movement} --- Contração muscular e função motora}

\begin{solution}{exo:b2:muscle-movement:1}
Os filamentos finos deslizam ao longo dos grossos, puxados pelas
cabeças de miosina; nenhum dos filamentos encurta. A banda I e a zona H
se estreitam e os discos Z se aproximam; a banda A (os filamentos
grossos) conserva seu comprimento.
\end{solution}

\begin{solution}{exo:b2:muscle-movement:2}
(1) ATP ligado, cabeça solta; ATP hidrolisado, cabeça armada. (2) A
cabeça se liga à actina. (3) Fosfato liberado, golpe de força, ADP
liberado. (4) Um novo ATP se liga e a cabeça se solta. O ATP é
necessário para \emph{soltar} a cabeça: sem ele, todas as cabeças
permanecem ligadas e o músculo trava --- é a rigidez.
\end{solution}

\begin{solution}{exo:b2:muscle-movement:3}
Impulso nervoso; liberação de acetilcolina e potencial de placa motora
(\qty{0.6}{ms}); \omterm{prop:b2:neurons-synapses:ap}{potencial de ação} muscular ao longo da fibra e pelos
\omterm{def:b2:cell-differentiation:fibre}{túbulos T} (\qty{1}{ms}); cálcio liberado pelo \omterm{def:b2:cell-differentiation:fibre}{retículo sarcoplasmático}
(\qty{1}{ms}); o cálcio se liga à troponina, a tropomiosina se desloca,
as cabeças se ligam (alguns milissegundos); a força sobe ao longo de
dezenas de milissegundos.
\end{solution}

\begin{solution}{exo:b2:muscle-movement:4}
\omterm{def:b2:muscle-movement:motorunit}{Unidade motora}: um \omterm{def:b2:neurons-synapses:neuron}{neurônio} motor e todas as fibras que ele inerva.
\omterm{def:b2:muscle-movement:motorunit}{Abalo}: a contração breve produzida por um único impulso. \omterm{def:b2:muscle-movement:motorunit}{Tétano}: a
contração fundida e sustentada produzida por impulsos que chegam mais
depressa do que o \omterm{def:b2:muscle-movement:motorunit}{abalo} decai. A força é graduada pela frequência de
disparo de cada unidade e pelo número de unidades recrutadas.
\end{solution}

\begin{solution}{exo:b2:muscle-movement:5}
Cada meio \omterm{def:b2:cell-differentiation:fibre}{sarcômero} desliza \qty{0.2}{\micro m} em \qty{50}{ms}:
\qty{4}{\micro m/s}; isso são 20 golpes de \qty{10}{nm} em
\qty{50}{ms}, 400 por segundo se uma cabeça ficasse ligada o tempo todo.
\end{solution}

\begin{solution}{exo:b2:muscle-movement:6}
\qty{2}{cm} em \qty{0.1}{s}: \qty{0.2}{m/s}, 2 comprimentos por segundo;
cada \omterm{def:b2:cell-differentiation:fibre}{sarcômero} encurta $2/\num{40000}$ cm $= \qty{0.5}{\micro m}$ em
\qty{0.1}{s}, \qty{5}{\micro m/s}.
\end{solution}

\begin{solution}{exo:b2:muscle-movement:7}
$F/F_0 = 0.3125/(v/v_{\max} + 0.25) - 0.25$: a 1, 3 e 6 comprimentos por
segundo ($0.1$, $0.3$, $0.6$ de $v_{\max}$): $0.64$, $0.32$, $0.12$.
Potência $Fv$ (em $F_0\times$ comprimentos por segundo): $0.64$, $0.95$,
$0.71$ --- máxima perto de 3 comprimentos por segundo, cerca de um terço
de $v_{\max}$.
\end{solution}

\begin{solution}{exo:b2:muscle-movement:8}
$80\times 30 = \qty{2400}{N}$; no pé, $2400\times 5/45 =
\qty{270}{N}$.
\end{solution}

\begin{solution}{exo:b2:muscle-movement:9}
A \qty{10}{Hz}, os impulsos chegam a cada \qty{100}{ms}, depois que o
cálcio do anterior desapareceu: \omterm{def:b2:muscle-movement:motorunit}{abalos} separados, cada um partindo de
um estado parcialmente relaxado, o que dá uma ondulação. A
\qty{50}{Hz}, eles chegam a cada \qty{20}{ms}, antes que o cálcio seja
bombeado: o cálcio permanece alto e a força se funde. O \omterm{def:b2:muscle-movement:motorunit}{tétano} supera o
\omterm{def:b2:muscle-movement:motorunit}{abalo} porque um único pulso de cálcio é breve demais para que as pontes
cruzadas estirem os elementos elásticos em série e desenvolvam a força
total; o cálcio sustentado o permite.
\end{solution}

\begin{solution}{exo:b2:muscle-movement:10}
Trabalho de \qty{20}{kJ}, energia do ATP de \qty{80}{kJ}, $1.6$ mol de
ATP. Primeiros \qty{5}{s}: $0.8$ mol de \omterm{prop:b2:muscle-movement:energy}{fosfocreatina}. Os $0.8$ mol de
ATP restantes pela glicólise: $0.4$ mol de glicose, $0.8$ mol de
lactato --- em \qty{40}{L}, \qty{20}{mmol/L}.
\end{solution}

\begin{solution}{exo:b2:muscle-movement:11}
Canal de liberação bloqueado: nenhum cálcio, nenhuma contração ---
paralisia. Bomba bloqueada: o cálcio permanece alto, o músculo não
consegue relaxar --- uma contratura. Troponina insensível: a
tropomiosina nunca se desloca, nenhuma contração. Sem cálcio no \omterm{def:b2:blood-circulation:blood}{sangue}:
o músculo esquelético se contrai normalmente (seu cálcio é interno); o
\omterm{def:b2:heart:anatomy}{coração}, que precisa de cálcio externo para desencadear sua liberação,
para.
\end{solution}

\begin{solution}{exo:b2:muscle-movement:12}
A cabeça converte até metade da energia livre de um ATP em trabalho; o
restante é calor, e as perdas das bombas, das pontes cruzadas que se
ligam sem puxar em alta velocidade e dos elementos elásticos levam o
conjunto a um quarto --- a eficiência de um bom motor a gasolina. Duas
maneiras de graduar a força permitem ao mesmo tempo um controle fino
em esforço baixo (pequenas unidades acrescentadas uma a uma, cada uma
disparando mais ou menos depressa) e uma grande amplitude (de mil
vezes, de uma pequena unidade na frequência do \omterm{def:b2:muscle-movement:motorunit}{abalo} a todas as unidades
em \omterm{def:b2:muscle-movement:motorunit}{tétano}); um único interruptor daria uma força ou nenhuma.
\end{solution}

\begin{solution}{pb:b2:muscle-movement:1}
\textbf{1.} $v = \sqrt{2\times 9.81\times 0.4} = \qty{2.8}{m/s}$;
$\tfrac12\times 70\times 2.8^{2} = \qty{275}{J}$.
\textbf{2.} Trabalho $= 275 + 70\times 9.81\times 0.4 = \qty{550}{J}$
ao longo de \qty{0.4}{m}: \qty{1370}{N}, o dobro do peso.
\textbf{3.} $F_0 = 500\times 30 = \qty{15000}{N}$; no pé,
\qty{1500}{N}.
\textbf{4.} \qty{4}{cm} em \qty{0.3}{s}: \qty{0.13}{m/s}, $1.1$
comprimento por segundo, $0.14\,v_{\max}$.
\textbf{5.} $F/F_0 = 0.3125/(0.14 + 0.25) - 0.25 = 0.55$: \qty{8300}{N}
no músculo, \qty{830}{N} no pé --- bem aquém dos \qty{1370}{N}
necessários. Só o encurtamento muscular não consegue realizar esse
salto. O contramovimento fornece o restante: quando quem salta se
agacha, os tendões estirados armazenam energia elástica, e o músculo,
contraindo-se de modo quase isométrico com alta força, os carrega; os
tendões então se retraem mais depressa do que qualquer fibra consegue
encurtar.
\textbf{6.} $550/0.3 = \qty{1.8}{kW}$; \qty{92}{W/kg} de extensor.
\textbf{7.} $4/\num{48000}$ cm $= \qty{0.83}{\micro m}$ por \omterm{def:b2:cell-differentiation:fibre}{sarcômero},
\qty{0.42}{\micro m} por meio \omterm{def:b2:cell-differentiation:fibre}{sarcômero}: 42 golpes de \qty{10}{nm}.
\textbf{8.} \qty{0.42}{\micro m} em \qty{0.3}{s}: \qty{1.4}{\micro m/s};
140 golpes por segundo se a cabeça ficasse ligada continuamente.
\textbf{9.} $500\times 6\times 10^{10}\times \num{48000}\times 300 =
4.3\times 10^{20}$ cabeças.
\textbf{10.} $0.55\times \num{15000}\times 0.04 = \qty{330}{J}$; por
golpe, $0.4\times 8\times 10^{-20} = \qty{3.2e-20}{J}$; $1.0\times
10^{22}$ golpes.
\textbf{11.} $1.0\times 10^{22}/4.3\times 10^{20} = 24$ golpes por
cabeça, dos 42 disponíveis: cerca de \qty{60}{\%}. A reserva é pequena
--- o músculo trabalha perto de seu limite, como constatou a questão 5.
\textbf{12.} $1.0\times 10^{22}/6.0\times 10^{23} = 0.017$ mol,
\qty{8.7}{g}; sua energia, $0.017\times 50 = \qty{860}{J}$, dos quais
\qty{330}{J} se tornaram trabalho: \qty{39}{\%}.
\textbf{13.} Mais rápido que um impulso a cada \qty{30}{ms}: acima de
\qty{33}{Hz}; cerca de 10 impulsos em \qty{0.3}{s}.
\textbf{14.} $9.9\times 10^{-6}\times 5\times 10^{-10}\times 6\times
10^{23} = 3\times 10^{9}$ íons livres; $1.5\times 10^{9}$ ATP.
\textbf{15.} Fibras: $500/(\pi\times(3\times 10^{-3})^{2}) = 1.8\times
10^{7}$; ATP das pontes cruzadas por fibra por impulso $1.0\times
10^{22}/(1.8\times 10^{7}\times 10) = 6\times 10^{13}$, quarenta mil
vezes a estimativa da bomba. O cálcio livre é um pequeno resto: o
retículo libera da ordem de \qty{0.1}{mmol/L}, dez vezes o aumento livre,
quase todo ele imediatamente ligado à troponina e aos tampões --- e todo
ele precisa ser bombeado de volta.
\textbf{16.} Para um salto máximo, quase todas as unidades em algumas
dezenas de milissegundos, as pequenas primeiro e as grandes por último;
um passo suave recruta talvez de um décimo a um quinto delas.
\textbf{17.} As unidades são recrutadas por ordem de tamanho: em esforço
baixo, cada nova unidade acrescenta um pequeno incremento; em esforço
alto, as unidades restantes são grandes e cada uma acrescenta um grande
degrau.
\textbf{18.} O potencial de placa motora de cada fibra cai à metade; as
fibras cujo potencial fica abaixo do limiar não disparam, de modo que o
\omterm{def:b2:muscle-movement:motorunit}{abalo} e o \omterm{def:b2:muscle-movement:motorunit}{tétano} são menores, e o salto não chega à altura.
\textbf{19.} $5\times 20 = \qty{0.1}{mol}$ de ATP: o equivalente a seis
saltos.
\textbf{20.} $20\times 20 = \qty{0.4}{mol}$ de \omterm{prop:b2:muscle-movement:energy}{fosfocreatina}: cerca de
23 saltos.
\textbf{21.} Dez saltos usam $0.17$ mol, ainda sem esgotar a
\omterm{prop:b2:muscle-movement:energy}{fosfocreatina}; depois de uns vinte, a glicólise assume, e o lactato e
os prótons se acumulam.
\textbf{22.} \qty{1.8}{kW} de trabalho a \qty{40}{\%} são \qty{4.6}{kW}
de ATP, $0.09$ mol por segundo, o que exige $0.015$ mol de oxigênio por
segundo --- \qty{0.34}{L/s}, \qty{21}{L/min}, dez vezes o que o \omterm{def:b2:blood-circulation:blood}{sangue}
consegue trazer: o salto é anaeróbio por necessidade, e o oxigênio vem
depois.
\textbf{23.} $0.17$ mol de ATP exigem $0.029$ mol de oxigênio,
\qty{0.64}{L}; pagos ao longo de um minuto, \qty{640}{mL/min}, duas
vezes e meia o consumo de repouso.
\textbf{24.} O ATP é refosforilado a partir da \omterm{prop:b2:muscle-movement:energy}{fosfocreatina} em
milissegundos, de modo que sua concentração quase não cai; a rigidez
exige um ATP perto de zero, o que um músculo vivo com alguma reserva
nunca atinge.
\textbf{25.} 42 golpes disponíveis por meio \omterm{def:b2:cell-differentiation:fibre}{sarcômero}, 24 necessários;
$F/F_0 = 0.55$ na velocidade do salto; $0.017$ mol, \qty{8.7}{g}, de
ATP.
\end{solution}
